\chapter{Sequence Models and Local Basis Methods}

Two important departures from standard feedforward learning appear in this part of the course. Recurrent neural networks handle temporal dependence by carrying forward a latent state. Radial basis function networks handle nonlinear regression and classification by stitching together local basis functions.

\section{Recurrent neural networks}

A basic recurrent network updates its hidden state according to
\[
  \hvec_t = \sigma(\Wmat_x \vec{x}_t + \Wmat_h \hvec_{t-1} + \vec{b}),
  \qquad
  \vec{y}_t = g(\Wmat_o \hvec_t + \vec{c}).
\]
The hidden state is a learned summary of the past. Unlike a fixed sliding window, it does not require the horizon of relevant history to be chosen in advance.

\begin{aside}[State as compressed memory]
The network does not store the whole past explicitly. It learns which aspects of the past are worth preserving in $\hvec_t$ for future predictions.
\end{aside}

\section{Backpropagation through time}

Training an RNN unfolds the recurrence across time and applies backpropagation to the resulting deep computational graph. If the sequence loss is
\[
  L = \sum_{t=1}^{T} \ell_t,
\]
then gradients at time $t$ depend on all later losses whose computation involved $\hvec_t$.

This is conceptually straightforward and numerically delicate. Repeated multiplication by the recurrent Jacobian can cause vanishing or exploding gradients, which makes long-range dependencies difficult to learn.

\section{Gating and stable memory}

Architectures such as LSTMs address this by introducing controlled memory paths. In a simplified form,
\[
  \vec{c}_t = \vec{f}_t \odot \vec{c}_{t-1} + \vec{i}_t \odot \tilde{\vec{c}}_t,
  \qquad
  \hvec_t = \vec{o}_t \odot \tanh(\vec{c}_t).
\]
The forget, input, and output gates regulate what is retained, written, and exposed. The mathematical role of gating is to create a route through time along which information can persist without being repeatedly distorted.

\section{Radial basis function networks}

An RBF network uses localized hidden units
\[
  \phi_i(\vec{x})
  =
  \exp\!\left(
    -\frac{\lVert \vec{x}-\vec{t}_i \rVert^2}{2\sigma_i^2}
  \right),
  \qquad
  y(\vec{x}) = \sum_{i=1}^{M} w_i \phi_i(\vec{x}).
\]
Each hidden unit is centered at a prototype $\vec{t}_i$ and responds only near that location.

\begin{keyequation}[Local approximation]
\[
y(\vec{x}) = \sum_{i=1}^{M} w_i \phi_i(\vec{x})
\]
\tcblower
RBF networks do not build a deep hierarchy. They tile the input space with local basis functions and then perform linear readout in that nonlinear feature space.
\end{keyequation}

The standard training strategy is two-stage:
\begin{enumerate}
  \item choose centers $\vec{t}_i$ by unsupervised clustering or prototypes,
  \item solve the output weights by linear regression once the basis functions are fixed.
\end{enumerate}

With design matrix $\Phi_{\alpha i} = \phi_i(\vec{x}^{(\alpha)})$, the least-squares solution satisfies
\[
  \Phi^\top \Phi \, \vec{w} = \Phi^\top \vec{y}.
\]

\section{What these models teach}

RNNs and RBF networks solve different problems, but they teach the same structural lesson: architecture is a form of prior knowledge. Recurrent connections encode temporal dependence. Radial basis functions encode locality. In both cases, performance depends on how well that prior matches the data geometry.
